\documentclass[12pt]{article}
\usepackage[utf8]{inputenc}
\usepackage[T1]{fontenc}
\usepackage[spanish]{babel}
\usepackage{amsmath}
\usepackage{amsfonts}
\usepackage{geometry}
\usepackage{xcolor}
\geometry{a4paper, margin=1in}

\title{\textbf{Guía Técnica AG-RACE}\\ Física de Hovercraft y Entrenamiento RL}
\author{Senior Physics \& RL Engineer}
\date{\today}

\begin{document}

\maketitle

\section{Introducción}
Este documento explica los fundamentos matemáticos del entorno AG-RACE, detallando el modelo de física tipo \textit{hovercraft} y el proceso de aprendizaje por neuroevolución utilizando el algoritmo CMA-ES con una arquitectura de red neuronal profunda.

\section{Resumen de Objetivos}
El objetivo principal del entrenamiento es crear un piloto que:
\begin{itemize}
    \item \textbf{Maximice la velocidad} sin perder el control (\textit{Drifting} controlado).
    \item \textbf{Mantenga la inercia}: Se penaliza el uso del freno para fomentar el flujo en curvas.
    \item \textbf{Sea agresivo} con otros competidores (el contacto físico es táctico).
    \item \textbf{Evite los muros}: El contacto con los límites resulta en eliminación inmediata.
\end{itemize}

\section{Física de Vuelo (Hover Physics)}
\subsection{Dinámica Longitudinal}
La velocidad de avance $v_f$ se calcula mediante:
\[ v_{f, t+1} = (v_{f, t} + A \cdot \tau \cdot \Delta t) \cdot \mu_{f} \cdot \delta \]
Donde $\tau$ es el acelerador y $\mu_{f}$ la fricción longitudinal.

\section{Entrenamiento Neuroevolutivo}
Utilizamos el algoritmo CMA-ES para optimizar una red neuronal profunda.

\subsection{Arquitectura Neuronal Gen 2 (Deep Scaling)}
La red \texttt{RacerPolicy} cuenta con **13,374 parámetros** distribuidos en 5 capas internas. Esta profundidad permite representar comportamientos más complejos.

\subsection{¿Hay Mutación y Reproducción?}
Aunque CMA-ES es un algoritmo estadístico, sus componentes mapean directamente a conceptos biológicos:

\begin{enumerate}
    \item \textbf{Reproducción (Muestreo)}: En lugar de "padre y madre", la reproducción ocurre mediante el muestreo de una distribución normal. Cada nuevo individuo es un "descendiente" de la media actual de la población.
    \item \textbf{Mutación (Sigma y Covarianza)}: 
    \begin{itemize}
        \item El parámetro $\sigma$ (step-size) controla la intensidad de la mutación global.
        \item La matriz de covarianza adapta la dirección de la mutación para "estirarse" hacia donde los resultados son mejores.
    \end{itemize}
    \item \textbf{Recombinación}: Al actualizar la media basándose en los mejores individuos, el algoritmo combina las mejores características de toda la generación.
\end{enumerate}

\subsection{Función de Fitness}
\[ F = (\text{Progreso} \times 2000) - (\text{Uso de Freno} \times 10) + (\text{Giro en Curva} \times 20) \]

\end{document}
